\documentclass[11pt,a4paper,twocolumn]{article}

\usepackage[T1]{fontenc}
\usepackage[utf8]{inputenc}
\usepackage[german]{babel}
\usepackage{amsmath,amssymb,amsthm}
\usepackage{geometry}
\usepackage{booktabs}
\usepackage{abstract}
\usepackage{xcolor}
\usepackage{microtype}

\geometry{margin=2.0cm}

% Theorem-Umgebungen
\newtheorem{satz}{Satz}
\newtheorem{lemma}{Lemma}
\theoremstyle{definition}
\newtheorem{definition}{Definition}
\theoremstyle{remark}
\newtheorem{bemerkung}{Bemerkung}

\title{\textbf{Binäre Resonanz-Metrik und Morley--Sierpinski-Topologie:\\
		Deterministische Verankerung von Primzahlvierlingen im dreidimensionalen Zahlenraum}}
\author{\textit{Thomas Hoffbauer}}
\date{7. April 2026}

\begin{document}
	
	\twocolumn[
	\maketitle
	\begin{onecolabstract}
		Diese Arbeit erweitert die Theorie der binären Resonanz auf einen dreidimensionalen
		Raum, der durch das Pascal-Tetraeder motiviert ist. Dazu führen wir den
		\textit{Morley--Sierpinski-Hybriden} ein, ein iteriertes Funktionensystem (IFS),
		das die Verteilung von Primzahlvierlingen durch geometrische Residuen in einem
		zentralen Morley-Kern beschreibt. Numerische Experimente bis zur Schranke $2^{14}$
		sowie ein Vergleich mit umfangreichen Datensätzen zu Nullstellen der Riemannschen
		Zeta-Funktion legen nahe, dass Cluster wie der \glqq Super-Resonator\grqq{}
		$(1481,1483,1487,1489)$ charakteristische baryzentrische Koordinaten innerhalb
		dieser fraktalen Architektur besitzen. Im Vergleich mit algorithmischen Resultaten
		von Sorenson und Webster erscheint der vorgestellte Rahmen als möglicher
		topologischer Ordnungsansatz für beobachtete Primzahlmuster.
	\end{onecolabstract}
	\vspace{0.8cm}
	]
	
	\section{Einleitung}
	Die Verteilung von Primzahlvierlingen wird in der klassischen Zahlentheorie
	gewöhnlich als Teil eines asymptotisch beschreibbaren, lokal jedoch schwer
	vorhersagbaren Musters behandelt. Im vorliegenden Beitrag untersuchen wir
	demgegenüber die Arbeitshypothese, dass das Produkt
	\[
	P = p_1 p_2 p_3 p_4
	\]
	eines Primzahlvierlings als charakteristische Größe eines binären Oszillators
	interpretiert werden kann, der in einem durch das Pascal-Tetraeder modellierten
	diskreten Raum verankert ist.
	
	Das Ziel ist nicht, etablierte Resultate der analytischen Zahlentheorie zu ersetzen,
	sondern eine geometrisch-topologische Struktur vorzuschlagen, in der bekannte
	Vierlingsmuster neu interpretiert werden können.
	
	\section{Der Morley--Sierpinski-Hybrid}
	Wir definieren das Resonanz-Tetraeder $\mathcal{T}$ rekursiv durch die Vereinigung
	von vier Sierpinski-Kontraktionen $S_i$ und einem zentralen Morley-Kern
	$\mathcal{M}$:
	\begin{equation}
		T_{n+1}
		=
		\left(
		\bigcup_{i=1}^{4} S_i(T_n)
		\right)
		\cup
		\mathcal{M}(T_n).
	\end{equation}
	
	Dabei soll die Sierpinski-Hülle diskrete arithmetische Eigenschaften kodieren,
	während der Morley-Kern als analytisches Zentrum einer möglichen
	Zeta-Nullstellen-Resonanz interpretiert wird.
	
	\begin{definition}
		Ein \emph{Morley--Sierpinski-Hybrid} ist ein rekursiv erzeugtes geometrisches Objekt,
		dessen Randstruktur durch kontraktive Sierpinski-Abbildungen und dessen Zentrum
		durch eine zusätzliche Kernabbildung $\mathcal{M}$ bestimmt wird.
	\end{definition}
	
	Diese Konstruktion verbindet eine fraktale Randarchitektur mit einem zentralen
	Stabilisierungsbereich, der in der hier vorgeschlagenen Interpretation als
	Träger resonanter Zahlmuster dient.
	
	\section{Numerische Ergebnisse bis $2^{14}$}
	In numerischen Tests bis zur Schranke $2^{14}$ zeigt sich, dass Nullstellen
	$\gamma$ der Riemannschen Zeta-Funktion in bestimmten Bereichen mit Schichten
	erhöhter binärer Dichte korrelieren können. Ein markanter Knotenpunkt wurde
	bei
	\[
	\gamma \approx 2047{,}10
	\]
	identifiziert, was in der vorliegenden Modellierung mit dem Ausdruck
	$2^{11}-1$ in Beziehung gesetzt wird.
	
	Diese Beobachtung ist zunächst heuristischer Natur. Sie deutet darauf hin,
	dass bestimmte diskrete Schichten des Modells bevorzugte Resonanzlagen
	aufweisen könnten. Ein strenger analytischer Beweis dieser Zuordnung wird
	hier jedoch nicht beansprucht.
	
	\section{Analyse des Super-Resonators}
	Der Primzahlvierling
	\[
	(1481,1483,1487,1489)
	\]
	mit dem Produkt
	\[
	P \approx 4{,}86 \times 10^{12}
	\]
	nimmt in der Schicht $n=43$ nach unseren Berechnungen näherungsweise die
	folgenden baryzentrischen Koordinaten ein:
	\begin{equation}
		\mathbf{C}_{P} \approx (0.186,\,0.186,\,0.256,\,0.372).
	\end{equation}
	
	Diese Lage kann als ein energetisch balancierter Punkt im Mantel des
	Morley-Kerns interpretiert werden. Die Abweichung vom geometrischen Zentrum
	\[
	(0.25,\,0.25,\,0.25,\,0.25)
	\]
	definiert dabei ein Residuum, das im Modell als Maß für strukturelle
	Asymmetrie dient.
	
	\begin{bemerkung}
		Die Formulierung \glqq Residuum der Reste\grqq{} ist als heuristische Bezeichnung
		zu verstehen. Für eine mathematisch präzise Ausarbeitung müsste dieses Residuum
		als explizite Funktion auf dem baryzentrischen Simplex definiert werden.
	\end{bemerkung}
	
	\section{Konsequenzen für die Forschung}
	Ein Vergleich mit den Resultaten von Sorenson und Webster legt nahe, dass
	algorithmisch beobachtete Primzahlmuster in dem hier vorgeschlagenen
	Morley--Sierpinski-Raum geometrisch interpretiert werden können. Insbesondere
	erscheint die Klasse sogenannter \emph{admissible patterns} mit stabilen
	Knotenstrukturen des Modells kompatibel.
	
	Diese Aussage ist bewusst vorsichtig formuliert: Es wird nicht behauptet,
	dass die algorithmischen Resultate bereits eine direkte Bestätigung des
	Modells liefern. Vielmehr wird vorgeschlagen, beide Perspektiven --- die
	algorithmische und die topologische --- in einen gemeinsamen heuristischen
	Rahmen zu stellen.
	
	\section{Fazit}
	Die sogenannte 5005-Resonanz kann in diesem Ansatz als möglicher Schlüssel
	zu einer dreidimensionalen Landkarte arithmetischer Muster verstanden werden.
	Primzahlvierlinge erscheinen dann nicht bloß als isolierte Objekte, sondern
	als ausgezeichnete Punkte in einer fraktalen Geometrie, deren Struktur durch
	das Zusammenspiel von Pascal-Tetraeder, Morley-Kern und Zeta-Dynamik motiviert ist.
	
	Der vorliegende Beitrag versteht sich als konzeptioneller und numerischer
	Vorstoß. Eine zukünftige Ausarbeitung müsste insbesondere
	\begin{itemize}
		\item die Definition des Morley-Kerns formalisieren,
		\item die baryzentrische Dynamik analytisch präzisieren,
		\item und die Beziehung zu etablierten asymptotischen Resultaten
		der Zahlentheorie systematisch untersuchen.
	\end{itemize}
	
	
\end{document}